% SPDX-FileCopyrightText: Copyright (c) 2026 NVIDIA CORPORATION & AFFILIATES. All rights reserved.
% SPDX-License-Identifier: Apache-2.0
%
% Report on geometry supervision for 3DGRUT. Built by scripts/report/build_report.sh, which
% generates every figure from the ablation run records first -- no number here is retyped from
% the write-up by hand except where a slide says so.
\documentclass[aspectratio=169,10pt]{beamer}

\usetheme{default}
\usecolortheme{seahorse}
\usepackage[T1]{fontenc}
\usepackage{helvet}
\renewcommand{\familydefault}{\sfdefault}
\usepackage{booktabs}
\usepackage{graphicx}
\usepackage{xcolor}
\usepackage{amsmath}

\setbeamertemplate{navigation symbols}{}
\setbeamertemplate{itemize items}[circle]
\setbeamercolor{frametitle}{fg=black,bg=gray!18}
\setbeamerfont{frametitle}{size=\large,series=\bfseries}
\setbeamerfont{footline}{size=\tiny}
\setbeamertemplate{footline}{%
  \hfill\insertframenumber\,/\,\inserttotalframenumber\hspace{2ex}\vspace{1ex}}

\definecolor{good}{HTML}{2E7D32}
\definecolor{bad}{HTML}{C62828}
\newcommand{\up}[1]{\textcolor{good}{#1}}
\newcommand{\dn}[1]{\textcolor{bad}{#1}}
\newcommand{\code}[1]{\texttt{\small #1}}

% Panel with a caption underneath, sized as a fraction of the text width.
\newcommand{\panel}[3]{%
  \begin{minipage}[t]{#1\textwidth}\centering
    \includegraphics[width=\linewidth]{#2}\\[0.2em]{\tiny #3}
  \end{minipage}}

% As \panel, but cropped: #2 is a trim spec in pixels of the source image, given as
% left bottom right top (the sources are 640x360, and 1px = 1bp here).
\newcommand{\panelcrop}[4]{%
  \begin{minipage}[t]{#1\textwidth}\centering
    \includegraphics[trim=#2, clip, width=\linewidth]{#3}\\[0.2em]{\tiny #4}
  \end{minipage}}

\title{\bfseries Geometry supervision for 3DGRUT}
\subtitle{Depth and normal quality on OB3D: what worked, what did not, and why}
\date{}
\author{Ablation report {\small — all numbers at 7k iterations on OB3D}}

\begin{document}

\begin{frame}[plain]
  \titlepage
  \vspace{-2em}
  \begin{center}\footnotesize
  Six candidate supervision terms. Two landed, two were refuted by measurement,\\
  and one refutation was of an idea introduced in this work.
  \end{center}
\end{frame}

\begin{frame}{The problem}
  3DGRUT reconstructs appearance well and geometry poorly. The rendered depth and normal
  buffers are byproducts of a purely photometric objective, so nothing constrains them to
  agree with the surface.

  \vspace{0.6em}
  \textbf{Setup for everything that follows}
  \begin{itemize}
    \item Three OB3D scenes with ground-truth depth and normals: \code{sponza},
          \code{lone-monk}, \code{emerald-square}.
    \item 7k iterations, gaussian primitive. Reference geometry is used \emph{only} for
          evaluation; no training term reads it.
    \item Seed noise, measured over \code{seed\_initialization} 1--4:
          $\pm$0.04--0.13\,dB PSNR, $\pm$0.002--0.004 \code{abs\_rel}, $\pm$0.1--1.2$^\circ$
          normal. Effects of a few tenths of a dB need repeated seeds.
  \end{itemize}

  \vspace{0.4em}
  \textbf{The normal metric needs a control.} A normal buffer is scored against the angular
  error of pointing every normal straight back along the view ray --- a ``normal'' that uses no
  geometry whatsoever. We report \code{n\_gain} $=$ control $-$ error. \textbf{Negative
  \code{n\_gain} means the buffer is worse than using no geometry at all.}
\end{frame}

\begin{frame}{Why the control matters: the baseline normal buffer is noise}
  \begin{center}
    \panel{0.3}{sponza/baseline_f4_rgb.png}{rendered RGB --- appearance is fine (36.3\,dB)}\hfill
    \panel{0.3}{sponza/baseline_f4_normal.png}{rendered normal --- \dn{high-frequency noise}}\hfill
    \panel{0.3}{sponza/baseline_f4_normal_gt.png}{reference normal}
  \end{center}
  \vspace{0.3em}
  Sponza, held-out view. The normal buffer is \emph{plausible at a glance} --- it has structure
  and colour variation --- and it scores 52.4$^\circ$ mean error with
  \code{n\_gain} $=$ \dn{$-9.6^\circ$}: worse than the geometry-free control.

  \vspace{0.4em}
  {\small This is what the whole effort is trying to fix, and it is why the raw angle is not
  reported alone. A term can improve the angle and still be losing to the control.}
\end{frame}

\begin{frame}{Methodology: the scenes fail differently, so scene averages lie}
  \begin{center}\small
  \begin{tabular}{lccl}
    \toprule
    Scene & floater px & \code{delta1} failures & Failure mode \\
    \midrule
    \code{sponza}         & 0.2--2\%  & moderate & floaters; per-ray spread detects them (AUC 0.68--0.90) \\
    \code{emerald-square} & up to 13\% & high    & floaters, large depth range \\
    \code{lone-monk}      & \textbf{0.01\%} & \textbf{17.7\%} & \textbf{confidently placed, wrong} --- no floaters to find \\
    \bottomrule
  \end{tabular}
  \end{center}

  \vspace{0.8em}
  \code{lone-monk} has essentially no floaters yet the worst \code{delta1} rate. Its depth is
  wrong in an \emph{opaque} way: the rays are already concentrated, just at the wrong distance.

  \vspace{0.6em}
  \textbf{Consequence, and a recurring theme of this report:} any term that works by
  concentrating a ray (``commit to one surface'') is structurally unable to help
  \code{lone-monk}. A scene average would have hidden that and reported a mild win.
\end{frame}

\begin{frame}{What was attempted}
  \begin{center}\small
  \begin{tabular}{llp{6.5cm}}
    \toprule
    Term & Outcome & One-line verdict \\
    \midrule
    Depth--normal consistency & \up{landed} & Transforms normals; best depth on \code{sponza} \\
    Depth variance along the ray & \dn{refuted} & Buys confidence, not accuracy \\
    Scale-$z$ regularisation (PGSR port) & \up{landed, small} & Only after fixing a wrong assumption \\
    Pseudo-depth \emph{regression} & \dn{abandoned} & Prior is worse than the model it would teach \\
    Pseudo-depth \emph{ordinal} & \up{landed} & Only term improving depth on all 3 scenes \\
    \quad + its confidence gate & \dn{refuted} & My own idea; measurement killed it \\
    Multi-view consistency, mesh export & not started & \\
    \bottomrule
  \end{tabular}
  \end{center}
  \vspace{0.5em}
  {\small Two of the five conclusions reverse a prediction made before measuring. Those are
  the slides worth reading.}
\end{frame}

\section{Landed: depth--normal consistency}

\begin{frame}{Depth--normal consistency: tie the normal to the depth gradient}
  Penalise disagreement between the rendered normal and the normal implied by the rendered
  depth's local geometry. Purely \emph{internal} --- no external data.
  \begin{center}
    \panel{0.3}{sponza/baseline_f4_normal.png}{baseline: 52.4$^\circ$, \code{n\_gain} \dn{$-9.6$}}\hfill
    \panel{0.3}{sponza/dn05_f4_normal.png}{\code{dn 0.05}: \up{25.4$^\circ$}, \code{n\_gain} \up{$+17.4$}}\hfill
    \panel{0.3}{sponza/baseline_f4_normal_gt.png}{reference}
  \end{center}
  \vspace{0.2em}
  The arch, the masonry courses and the boss are all recovered from a buffer that was noise.
  Depth improves too: \code{abs\_rel} \up{$-17.0\%$} on \code{sponza}.

  \vspace{0.3em}
  {\small Cost: $-0.2$\,dB PSNR. And note the scene: on \code{lone-monk} the same term gives
  \dn{$-0.5\%$} depth --- nothing. Section on compounding returns to this.}
\end{frame}

\section{Refuted: depth variance along the ray}

\begin{frame}{Depth variance: a sharpening prior with no reference for \emph{where}}
  Penalising per-ray depth variance asks each ray to commit to a single surface. It looks like
  a geometry term and behaves like a confidence term.

  \vspace{0.5em}
  \begin{columns}[T]
    \begin{column}{0.52\textwidth}
      \textbf{Why it cannot work as posed}
      \begin{itemize}\small
        \item Variance is zero for a Dirac at \emph{any} distance, so the objective says commit,
              never where.
        \item In pairwise form,
              $M_2 - D^2/\!\mathrm{acc} = \frac{1}{2\,\mathrm{acc}}\sum_{ij} w_i w_j (t_i-t_j)^2$,
              the two-hit case is a \textbf{double well} in the near opacity: $w_\mathrm{near}
              w_\mathrm{far}$ carries an $\alpha(1-\alpha)$ factor peaking at $\tfrac12$.
        \item The near well is \textbf{absorbing}: once the front particle reaches $\alpha=1$,
              transmittance zeroes the gradient to everything behind it, for every loss. The
              error is permanent.
      \end{itemize}
    \end{column}
    \begin{column}{0.45\textwidth}
      \textbf{Measured consequence}
      \begin{center}\small
      \begin{tabular}{lc}
        \toprule
        & change \\
        \midrule
        per-ray spread & $\times$\,0.17 (6$\times$ tighter) \\
        \code{wrong\,|\,tight} & \dn{$\times$\,170} \\
        \code{sponza} \code{abs\_rel} & \dn{$+6\%$} \\
        floater fraction & \dn{0.30\% $\to$ 1.62\%} \\
        \bottomrule
      \end{tabular}
      \end{center}
      \vspace{0.4em}
      {\small It did exactly what it was asked --- rays six times tighter --- and became 170
      times more likely to be confidently wrong.}
    \end{column}
  \end{columns}
\end{frame}

\begin{frame}{The same term, rescaled: minima unchanged, behaviour changed}
  Dividing by $\mu^2$ leaves \emph{both} wells at exactly zero. That was first dismissed as a
  non-fix --- same minima, same problem.

  \vspace{0.5em}
  It moved the \textbf{separatrix} from 0.49 to 0.82, shrinking the bad basin from 51\% of the
  axis to 18\%. Descent responds to the barrier, not to the location of the optima.
  \code{sponza} goes \dn{$+6\%$} $\to$ \up{$-7.8\%$}.

  \vspace{0.8em}
  \textbf{Two lessons that outlived the term}
  \begin{itemize}
    \item ``The minima are unchanged'' is not an argument that behaviour is unchanged. When
          judging a reparameterisation of a degenerate objective, locate the separatrix.
    \item $\mathrm{Var}/\mu^2 = \mathrm{acc}\cdot M_2/D^2 - 1$ needed \emph{no} new buffer,
          kernel or backward pass, while the mip-NeRF/2DGS kernel it replaced would have needed
          a new accumulator on three separate backward paths to reach a strictly worse place.
          Check what the existing buffers already express before extending the renderer.
  \end{itemize}
\end{frame}

\section{Landed: ordinal pseudo-depth}

\begin{frame}{The plan called for a scale-invariant regression loss. Measurement said no.}
  \begin{columns}[T]
    \begin{column}{0.46\textwidth}
      \begin{center}
        \includegraphics[width=\linewidth]{plots/prior_alignment.pdf}
      \end{center}
    \end{column}
    \begin{column}{0.52\textwidth}
      \code{DepthAnythingV2-Base} on \code{sponza}, affine-aligned against ground truth ---
      the \emph{most favourable} alignment possible, since it uses the ground truth the loss
      would not have.

      \vspace{0.5em}
      \begin{itemize}\small
        \item One affine per frame: \dn{0.068}, \textbf{worse than the 0.058 model it would be
              teaching}. An aligned L1 would have spent a foundation-model dependency to make
              depth worse.
        \item Per 16$\times$16 patch: \up{0.0116}. The local structure is excellent; the global
              structure drifts.
        \item So the signal is not in \emph{any} single affine gauge, and scale-invariance ---
              the exact remedy the plan specified --- would not have saved it.
      \end{itemize}
    \end{column}
  \end{columns}
  \vspace{0.4em}
  {\small Two further traps: the model emits \emph{disparity} (monotonically inverted, not
  merely mis-scaled) and \emph{z-depth}, while this renderer's convention is Euclidean ray
  distance (20\% divergence at the corners).}
\end{frame}

\begin{frame}{What landed instead: read only the ordering}
  For a pair of pixels, ask whether the render agrees with the prior about \emph{which is
  nearer}, and penalise the rendered gap only where they disagree:
  \[
    \mathcal{L} = \frac{1}{|\mathcal{P}|}\sum_{(a,b)\in\mathcal{P}}
      \mathrm{relu}\!\big(-(d_a - d_b)\cdot \mathrm{sign}(-(\mathrm{disp}_a - \mathrm{disp}_b))\big)
      \;/\; \text{scene extent}
  \]
  \begin{itemize}
    \item \textbf{Invariant to any increasing transform} of the prior, not just an affine one.
          The alignment problem, the disparity inversion (one sign flip) and the
          z-vs-ray-distance mismatch all disappear at once.
    \item \textbf{One-sided}: a pair the render already orders correctly contributes exactly
          zero \emph{and no gradient}, so the term cannot fight geometry it agrees with.
    \item Pairs are formed by \emph{cropping} at a random offset $\leq$5\% of the image side,
          not by \code{torch.roll} --- wrapping pairs opposite edges and manufactures
          disagreements the prior never claimed.
  \end{itemize}
  \vspace{0.3em}
  {\small On the pixels where the trained model fails \code{delta1}, the prior is closer to
  ground truth \textbf{91\%} of the time --- the population this targets is real.}
\end{frame}

\begin{frame}{Result: the first term to improve depth on all three scenes}
  \begin{center}\small
  \begin{tabular}{lccc}
    \toprule
    \code{abs\_rel}, mean $\pm$ sd over 3 seeds & \code{sponza} & \code{lone-monk} & \code{emerald-square} \\
    \midrule
    baseline            & 0.0571 $\pm$ 0.0009 & 0.0982 $\pm$ 0.0019 & 0.1368 $\pm$ 0.0010 \\
    ordinal, $\lambda=0.1$ & \up{0.0506 $\pm$ 0.0008} & \up{0.0876 $\pm$ 0.0010} & \up{0.1228 $\pm$ 0.0033} \\
    & \up{$-11.4\%$} & \up{$-10.8\%$} & \up{$-10.3\%$} \\
    \midrule
    PSNR change & \up{$+0.12$\,dB} & \dn{$-0.20$\,dB} & \dn{$-0.61$\,dB} \\
    \bottomrule
  \end{tabular}
  \end{center}
  \vspace{0.3em}
  A consistent 10--11\% at 4--8$\sigma$. Every earlier term had to be argued scene by scene.
  \begin{center}
    \includegraphics[width=0.78\textwidth]{plots/lambda_sweep.pdf}
  \end{center}
  \vspace{-0.4em}
  {\small The usable band is narrow: $\lambda=1$ buys more depth on \code{lone-monk} and costs
  \code{emerald} 2.4\,dB; $\lambda=10$ gives up the gain entirely.}
\end{frame}

\begin{frame}{Qualitatively: what the 10\% is, and what it is not}
  \begin{center}
    \panel{0.24}{lone-monk/gt_f4_depth.png}{reference depth}\hfill
    \panel{0.24}{lone-monk/baseline_f4_depth.png}{baseline}\hfill
    \panel{0.24}{lone-monk/baseline_f4_depth_err.png}{baseline error}\hfill
    \panel{0.24}{lone-monk/pd01_f4_depth_err.png}{ordinal error}
  \end{center}
  \vspace{0.4em}
  \code{lone-monk}. Reference depth shows the arch openings as \emph{far} (orange); the
  baseline puts a confident surface across them at wall depth (green). That is the
  ``confidently placed, wrong'' failure no ray-concentration term can reach.

  \vspace{0.5em}
  \textbf{Honest reading of the error maps:} the arch openings stay \dn{saturated} --- the
  ordinal term does \emph{not} fix them. What improves is the diffuse error everywhere else:
  the roof line, the walls and the ground all go blue.

  \vspace{0.3em}
  {\small That is the mechanism working as measured, not a disappointment. Pairs are drawn at
  $\leq$5\% offsets and the prior is reliable \emph{locally} (0.0116 per patch), so the term
  corrects regional depth. A hole in a wall is a long-range discontinuity --- exactly the
  regime where the prior drifts.}
\end{frame}

\section{Refuted: my own gate}

\begin{frame}{The gate: an idea of mine that measurement killed}
  \begin{columns}[T]
    \begin{column}{0.44\textwidth}
      \textbf{The reasoning, before training}\\[0.3em]
      {\small Near-tied pairs are where the prior's sign is noise. Drop pairs whose disparity
      gap is below 0.05 of the map's IQR:
      \begin{itemize}
        \item ordinal agreement with GT: 84\% $\to$ \textbf{97\%}
        \item so ungated, one pair in six pushes the wrong way
      \end{itemize}
      This looked decisive. It is wrong.}
    \end{column}
    \begin{column}{0.54\textwidth}
      \begin{center}
        \includegraphics[width=\linewidth]{plots/gate.pdf}
      \end{center}
    \end{column}
  \end{columns}
  \vspace{0.3em}
  \textbf{Why it fails}, both visible in numbers already collected:
  \begin{itemize}\small
    \item A large disparity gap selects a \emph{long-range} comparison --- precisely where this
          prior drifts (0.068 global vs 0.0116 per patch). The gate keeps the prior's weakest
          structure and discards its best.
    \item For a \emph{one-sided} loss, 97\% agreement is a warning, not a benefit: only 3\% of
          surviving pairs can produce any gradient at all.
  \end{itemize}
  {\small The offline metric counted pairs instead of asking what they teach. Default is now
  off, pinned by a test so re-enabling it has to be re-measured.}
\end{frame}

\section{Compounding}

\begin{frame}{The terms fix different scenes; they do not add}
  \begin{columns}[T]
    \begin{column}{0.44\textwidth}
      \begin{center}
        \includegraphics[width=\linewidth]{plots/scene_split.pdf}
      \end{center}
    \end{column}
    \begin{column}{0.54\textwidth}
      \textbf{Internal consistency vs external prior}
      \begin{itemize}\small
        \item Depth--normal consistency owns \code{sponza} (\up{$-17.0\%$}) and is
              \dn{null} on \code{lone-monk} (\dn{$-0.5\%$}).
        \item The ordinal prior is the reverse: even on \code{lone-monk} it gives
              \up{$-10.8\%$}.
      \end{itemize}
      \vspace{0.5em}
      This is the scene-split slide made concrete. \code{lone-monk}'s error is not an
      inconsistency, so no internal condition detects it; it takes information from outside the
      scene.
      \vspace{0.5em}
      \textbf{Combined, depth is sub-additive} --- the pair lands near the \emph{max} of the
      two, never the sum, and on \code{lone-monk} it is \dn{worse than the prior alone}
      ($-7.1\%$ vs $-10.8\%$).
    \end{column}
  \end{columns}
\end{frame}

\begin{frame}{But on normals the pairing is where everything pays off}
  \begin{center}
    \includegraphics[width=0.92\textwidth]{plots/compounding.pdf}
  \end{center}
  \vspace{-0.2em}
  The ordinal term \textbf{alone does nothing for normals} (52.4$^\circ\to$52.5$^\circ$, still
  below the control) --- expected, since the rendered normal is the particle's local $z$ axis
  and depth ordering does not reach it.

  \vspace{0.3em}
  \textbf{Yet \code{pd + dn} is the only configuration whose normals clear the control on all
  three scenes.} Depth--normal consistency ties the normal to the depth \emph{gradient}, so it
  can only be as good as the depth handed to it: the prior improves the depth, and consistency
  carries that across into normals.
\end{frame}

\begin{frame}{The cross-channel effect, seen directly}
  \begin{center}
    \panel{0.3}{lone-monk/dn05_f4_normal.png}{\code{dn} alone: 37.4$^\circ$, \code{n\_gain} \dn{$-2.1$}}\hfill
    \panel{0.3}{lone-monk/pd01_dn_f4_normal.png}{\code{pd + dn}: \up{28.3$^\circ$}, \code{n\_gain} \up{$+6.9$}}\hfill
    \panel{0.3}{lone-monk/baseline_f4_normal_gt.png}{reference}
  \end{center}
  \vspace{0.4em}
  \code{lone-monk} normals. With depth--normal consistency alone the columns and arches are
  blobby and the buffer \textbf{still loses to the geometry-free control}
  (\code{n\_gain} \dn{$-2.1^\circ$}). Adding a depth term that touches no normals resolves the
  colonnade and flips the sign (\up{$+6.9^\circ$}).

  \vspace{0.5em}
  {\small The lesson generalises: before concluding a depth term is irrelevant to normals,
  check the cross-channel pairing.}
\end{frame}

\begin{frame}{A methodological trap worth one slide}
  The pairing of the ordinal term with relative depth variance was first measured at
  $\lambda=1$ for \textbf{both} --- 10$\times$ and 100$\times$ their measured optima --- and
  written up as a null.

  \vspace{0.6em}
  \begin{center}\small
  \begin{tabular}{lccc}
    \toprule
    depth \code{abs\_rel} change & \code{sponza} & \code{lone-monk} & \code{emerald-square} \\
    \midrule
    ordinal alone ($\lambda=0.1$)        & $-11.4\%$ & $-10.8\%$ & $\mathbf{-10.3\%}$ \\
    variance alone ($\lambda=0.01$)      & $-7.9\%$ & $+1.2\%$ & $-3.2\%$ \\
    both, at their own optima            & \up{$\mathbf{-13.9\%}$} & $-10.4\%$ & \dn{$-6.8\%$} \\
    both at $\lambda=1$ (the first pass) & \dn{$-7.6\%$} & $-12.9\%$ & \dn{$+10.3\%$} \\
    \bottomrule
  \end{tabular}
  \end{center}

  \vspace{0.5em}
  At sane weights the verdict changes from \emph{harmful} to \emph{scene-dependent}: the pair
  beats both single terms on \code{sponza}, is a wash on \code{lone-monk}, and is still clearly
  \dn{worse than the prior alone} on \code{emerald-square}. That is a weaker claim than the
  first pass suggested in the other direction, and it is the one the data supports.

  \vspace{0.3em}
  \textbf{A null from a combination measured at the wrong weights is not a null about the
  combination.} The over-weighted cell is kept in the sweep so the trap stays visible.
\end{frame}

\begin{frame}{The photometric cost, seen directly: it is a real regression}
  \begin{center}
    % The band with the largest error increase, located by differencing both renders against
    % the reference rather than by eye: x in [0,240], y in [60,240] of the 640x360 frame.
    \panelcrop{0.31}{0 120 400 60}{emerald-square/gt_f4_rgb.png}{reference}\hfill
    \panelcrop{0.31}{0 120 400 60}{emerald-square/baseline_f4_rgb.png}{baseline --- 32.9\,dB}\hfill
    \panelcrop{0.31}{0 120 400 60}{emerald-square/pd01_dn_f4_rgb.png}{\code{pd + dn} --- \dn{32.0\,dB}}
  \end{center}
  \vspace{0.6em}
  \code{emerald-square} foliage, the scene and the region carrying the largest cost. Two things
  happen at once:
  \begin{itemize}\small
    \item The baseline's \textbf{spiky starburst artifacts} (left of frame, mid-left) are
          \up{gone} --- they are not in the reference either.
    \item The genuine foliage structure is \dn{smoothed away} with them.
  \end{itemize}
  \vspace{0.3em}
  \textbf{Against the reference the blur costs more than the artifacts did}: mean absolute
  error 4.28 $\to$ \dn{4.80}, with \code{pd + dn} closer to the reference on only \dn{41\%} of
  pixels. Vegetation is where a consistency term assuming a locally planar surface should be
  expected to pay most, and it does.
\end{frame}

\begin{frame}{Recommended settings}
  \begin{center}\small
  \begin{tabular}{lll}
    \toprule
    Goal & Configuration & Cost \\
    \midrule
    Depth, one term, no tuning & \code{pd 0.1} alone & 0--0.6\,dB \\
               & $-10$ to $-11\%$ on all three scenes; the safe default & \\
    \addlinespace
    Depth \emph{and} normals & \code{pd 0.1} $+$ \code{dn 0.05} & \dn{0.2--0.9\,dB} \\
               & the only config clearing the normal control on all 3 scenes, & \\
               & and the best \code{sponza} depth measured (\up{$-17.3\%$}) & \\
    \bottomrule
  \end{tabular}
  \end{center}
  \vspace{0.5em}
  \code{pd 0.1} $+$ \code{dvrel 0.01} is the best \emph{depth-only} pairing on \code{sponza}
  (\up{$-13.9\%$}) at a smaller photometric cost (0.1--0.2\,dB), but it loses to \code{pd}
  alone on \code{emerald-square}, so it is not a general recommendation.

  \vspace{0.4em}
  All of these remain \textbf{off by default}. The photometric cost is real, and PSNR is what
  the codebase is currently optimised and evaluated for.
\end{frame}

\begin{frame}{Limitations, stated plainly}
  \begin{itemize}
    \item \textbf{Everything is at 7k iterations.} No term has been confirmed at 30k. Gates
          tuned for a 7k sweep (\code{*\_from\_iter}) do not transfer unchanged.
    \item \textbf{Three scenes.} They were chosen because they fail differently, not because
          they are representative.
    \item \textbf{Combination cells are single-seed}; single-term cells are 3-seed. The
          \code{lone-monk} normal sign flip and the \code{sponza} depth gains are far outside
          seed noise; \code{emerald} differences of $\sim$1 point are not.
    \item \textbf{The emerald-square PSNR cost (0.6--0.9\,dB) is only half explained.} It is
          spatially concentrated in foliage, which is measured. \emph{Why} that scene suffers
          most --- largest depth range, so the prior's global drift should hurt most --- remains
          a hypothesis.
    \item \textbf{No mesh.} Nothing here has been evaluated by surface extraction or Chamfer
          distance, which is what ``better geometry'' would ultimately mean.
  \end{itemize}
\end{frame}

\begin{frame}{What the exercise actually taught}
  \begin{enumerate}
    \item \textbf{Always include a control that uses no information.} The baseline normal
          buffer looked reasonable and was worse than pointing every normal at the camera. Half
          the terms here would have been ``improvements'' without \code{n\_gain}.
    \item \textbf{Diagnose the scene, not the average.} \code{lone-monk} has no floaters and
          the worst depth; a whole family of terms is structurally incapable of helping it, and
          a scene mean would have reported a win.
    \item \textbf{An offline metric that motivates a mechanism can be measuring the wrong
          thing.} The gate's 84\%$\to$97\% was real and irrelevant: it counted pairs instead of
          asking what they teach.
    \item \textbf{Check what the reference quantity means \emph{here}.} ``Monocular depth'' is
          not depth up to one affine; PGSR's ``min scale'' is not this renderer's normal axis.
          Both cost a regression before being caught.
    \item \textbf{Measure combinations at each term's own optimum}, and against each single
          term --- not only against the baseline.
  \end{enumerate}
\end{frame}

\begin{frame}[plain]
  \vfill
  \begin{center}
    {\large\bfseries Summary}\\[1em]
    \begin{minipage}{0.86\textwidth}\small
    Two terms landed. Ordinal supervision from a monocular prior improves depth
    10--11\% on all three scenes --- the only term here to help all of them --- and, paired with
    depth--normal consistency, produces the only normal buffer that beats a geometry-free
    control everywhere.

    \vspace{0.8em}
    Three ideas were refuted by measurement, one of them introduced in this work. The refutations
    are recorded in \code{docs/normal-supervision.md} alongside the predictions they contradict,
    because the reasoning that failed is the part that generalises.
    \end{minipage}
  \end{center}
  \vfill
\end{frame}

\end{document}
